We adopt the standard agentic-LLM adversary model from Greshake et al.
{[}@greshake2023not{]} and Debenedetti et al.
{[}@debenedetti2024agentdojo{]}, with explicit cryptographic
assumptions.

\textbf{Adversary capabilities.} The adversary controls any untrusted
input the agent ingests: the user message, the contents of RAG-retrieved
documents, the output of any third-party tool whose result is fed back
into the agent's context, web pages the agent browses, and email or chat
messages the agent reads. The adversary observes the agent's system
prompt, the list of tools available to the agent, the JSON Schema of
every tool, and --- in our most pessimistic configuration --- the
constraint policy carried by the capability token in force at the time
of the attack. The adversary cannot forge Ed25519 signatures, cannot
compromise the issuer's private signing key, cannot execute code inside
the gate process, and cannot tamper with the audit log after the fact.

\textbf{Defender capabilities.} The defender controls a gate process
that sits between the agent runtime and the tool implementations. The
gate runs in a separate trust domain from the model --- a separate
process, container, or address space --- and is the only path through
which tools are invoked. The defender holds the issuer's private key in
a hardware security module or equivalent. The defender authors tool JSON
Schemas and constraint policies; we make no assumption that either is
correct, only that the SMT prover catches the cases where they disagree.

\textbf{Token confidentiality.} Capability tokens are \emph{unforgeable}
but not \emph{confidential}: an attacker who steals a valid token can
use it within its bounds. Production deployments should bind each token
to a session via a confidential channel (mTLS between agent and gate;
per-session token rotation; short TTLs) so that exfiltrated tokens are
useless to attackers operating in different sessions. The reference
implementation does not enforce this binding --- it is the
responsibility of the deployment. The threat model in this paper assumes
tokens cannot be exfiltrated from the agent runtime; this assumption is
realistic for in-process deployments and weakens to ``tokens have short
TTL and are bound to session keys'' in the multi-tenant case.

\textbf{Security goal.} For every tool call
\(c = (\text{tool}, \text{args})\) that the gate authorises with token
\(t\), the arguments \(\text{args}\) satisfy every constraint in \(t\),
and \(t\) is a valid descendant of a root token signed by a pinned
issuer. We refer to this property as \emph{gate soundness}. We also
require \emph{attenuation soundness}: no child token derived from a
parent can carry permissions broader than the parent along any
constraint dimension or lifetime.

\textbf{Operational definition of ``tool-call-mediated.''} Throughout
this paper an attack is \emph{tool-call-mediated} iff its success
criterion is the invocation of a specific tool with specific argument
values. This matches AgentDojo's \texttt{security} boolean and
InjecAgent's success predicate exactly. An attack whose success
criterion is, for example, the agent \emph{speaking aloud} a secret it
previously retrieved through a tool --- without subsequently emitting a
structured tool call to exfiltrate it --- is \textbf{not}
tool-call-mediated under this definition, regardless of whether a tool
call earlier in the trace produced the secret. We make no claim against
the non-tool-call-mediated class; see §6.5.

\textbf{Out of scope.} We list explicitly what VCD does \emph{not}
claim. Free-form natural-language output attacks --- for example, an
attacker coercing the agent to reveal a secret in its visible response
--- remain unaddressed; the bounded-grammar argument does not apply.
Side-channel exfiltration through allowed parameter values (an
\texttt{amount} field accepting any 64-bit integer permits up to 64 bits
per call of covert content) is not prevented; we discuss in Section 8
how schema tightening reduces this channel. Confused-deputy attacks
across multiple legitimately-authorised tool calls --- read mailbox,
then exfiltrate via legitimate \texttt{send\_email} --- are partially
addressed by our companion provenance layer {[}omitted for blind
review{]} but are not the focus of this paper. The model itself is not
part of the trusted computing base; the gate is. Compromise of the gate
process or the issuer's signing key is outside the model.

\textbf{Trust boundary of the gate.} Because the gate's integrity is
load-bearing for every claim in this paper, we are explicit about what
its operational profile must look like. The gate runs in a separate
process from the agent runtime, with no \texttt{eval}, \texttt{exec}, or
untrusted-deserialisation entry points reachable from agent-controlled
data; the issuer's private key resides in a hardware security module or
analogous trusted environment; the audit log is written to append-only
storage. Standard process-isolation, language-runtime hardening, and
key-management practices apply. These requirements are the same as for
any cryptographic gateway and are not specific to VCD. An agent runtime
compromised by an unrelated vulnerability --- RCE into the gate's
address space, key exfiltration via memory disclosure --- would defeat
the property; we are not protecting against operating-system-level
compromise.

This threat model is deliberately narrower than the universal ``we
defend against all prompt injection'' claims common in the prior
literature. We argue that narrow, well-defined threat models with formal
guarantees are more useful to deployers than broad informal claims with
statistical guarantees.
